\documentclass[11pt]{article}
\usepackage[margin=1in]{geometry}
\usepackage{amsmath,amssymb,amsthm,mathtools,esint}
\usepackage[hidelinks]{hyperref}
\usepackage{microtype}

\newtheorem{theorem}{Theorem}
\newtheorem{lemma}[theorem]{Lemma}
\newtheorem{remark}[theorem]{Remark}
\newcommand{\avg}{\mathop{\fint}}
\newcommand{\D}{\mathcal D}
\newcommand{\M}{M^{\mathrm d}_{Q_0}}
\newcommand{\JN}{JN^{\mathrm d}_p(Q_0)}
\newcommand{\Aone}{A^{\mathrm d}_1(Q_0)}
\newcommand{\Rop}{\mathcal R}

\title{A Coifman--Rochberg-Type Characterization of Dyadic
John--Nirenberg Space}
\author{Autonomous mathematical discovery run \texttt{fa\_banach\_001}}
\date{13 August 2026}

\begin{document}
\maketitle

\begin{abstract}
Kinnunen and Myyryl\"ainen ask whether dyadic John--Nirenberg space admits a
Coifman--Rochberg-type characterization.  We prove that, on every finite cube
and for $1<p<\infty$, dyadic $JN_p$ is exactly the difference of its positive
dyadic $A_1$ cone.  The decomposition is quantitative and its $A_1$ constants
are uniform.  Each positive summand also has an exact representation as a
bounded positive factor times a fractional power of a dyadic maximal
function.  The proof is a Rubio de Francia iteration made possible by the
maximal-operator theorem in the source paper and the weak-$L^p$ embedding of
$JN_p$.
\end{abstract}

\section{The question and the answer}

Let $Q_0\subset\mathbb R^n$ be a finite cube, let $\D(Q_0)$ be its dyadic
subcubes, and let $1<p<\infty$.  Write
\[
 [f]_{\JN}:=
 \sup_{\mathcal F}
 \left(\sum_{Q\in\mathcal F}|Q|
   \left(\avg_Q|f-f_Q|\right)^p\right)^{1/p},
\]
where the supremum is over pairwise disjoint families
$\mathcal F\subset\D(Q_0)$.  Constants are the kernel of this seminorm.  We
use the normalized Banach norm
\begin{equation}\label{eq:Xnorm}
 \|f\|_X:=|Q_0|^{1/p}|f_{Q_0}|+[f]_{\JN}.
\end{equation}
For a nonnegative locally integrable $w$, write $w\in\Aone$ when
\[
 \M w\leq [w]_{\Aone}w\qquad\text{a.e. on }Q_0.
\]

On page 2 of arXiv:2107.00492v2, the source paper asks whether a
Coifman--Rochberg-type characterization exists for dyadic John--Nirenberg
space.  It proves that $\M$ is bounded on the $JN_p^{\mathrm d}$ seminorm
and, motivated by the inclusion
$(\M h)^{1/p}\in JN_p^{\mathrm d}$ for $h\in L^1$, asks whether a
Coifman--Rochberg-type characterization exists.  The following gives a
precise affirmative answer.

\begin{theorem}[Dyadic $A_1$ difference characterization]\label{thm:main}
For every finite cube $Q_0$ and $1<p<\infty$,
\begin{equation}\label{eq:characterization}
 \JN=(\JN\cap\Aone)-(\JN\cap\Aone).
\end{equation}
More quantitatively, there is $C=C(n,p)$ such that every real
$f\in\JN$ can be written $f=u-v$, where $u,v>0$ a.e.,
\[
 u,v\in\JN\cap\Aone,
 \qquad [u]_{\Aone}+[v]_{\Aone}\leq C,
 \qquad \|u\|_X+\|v\|_X\leq C\|f\|_X.
\]
Conversely, every such difference belongs to $\JN$, and
\[
 \|f\|_X\leq \|u\|_X+\|v\|_X.
\]

There are also an exponent $r=r(n,p)>1$, functions
$h_+,h_-\in L^1_+(Q_0)$, and measurable multipliers $k_+,k_-$ satisfying
$C^{-1}\leq k_\pm\leq1$ such that
\begin{equation}\label{eq:maximalrepresentation}
 f=k_+(\M h_+)^{1/r}-k_-(\M h_-)^{1/r},
\end{equation}
and the two displayed summands are precisely $u$ and $v$ above.
\end{theorem}

\section{Idea of the proof}

The source theorem controls the oscillation seminorm of $\M f$, but the
Rubio de Francia algorithm needs a genuine Banach norm.  The weak-$L^p$
embedding of $JN_p$ supplies the missing average estimate, so $\M$ is bounded
for \eqref{eq:Xnorm}.  Absolute value is also bounded for this norm.

Starting from $g=|f|$, iterate the dyadic maximal operator:
\[
 \Rop g=\sum_{j\geq0}\frac{(\M)^j g}{(2A)^j},
 \qquad A=\|\M\|_{X\to X}.
\]
Then $\Rop g\geq g$, its $X$ norm is controlled, and
$\M(\Rop g)\leq2A\Rop g$.  Thus $\Rop|f|$ is an $A_1$ majorant of $|f|$.
The two functions $2\Rop|f|+f$ and $2\Rop|f|$ are positive, remain in
$JN_p$, and are comparable to that majorant, hence are uniformly $A_1$.
Finally, reverse H\"older self-improvement of $A_1$ weights converts each
summand to the maximal-power form \eqref{eq:maximalrepresentation}.

\section{The two norm estimates}

\begin{lemma}[Absolute value]\label{lem:absolute}
There is an absolute constant $C$ such that
$\||f|\|_X\leq C\|f\|_X$.
\end{lemma}

\begin{proof}
For every dyadic $Q$ and every constant $c$,
\[
 \avg_Q\bigl||f|-(|f|)_Q\bigr|
 \leq2\avg_Q\bigl||f|-|c|\bigr|.
\]
Taking $c=f_Q$ gives
$[|f|]_{\JN}\leq2[f]_{\JN}$.  On $Q_0$,
\[
 |Q_0|^{1/p}(|f|)_{Q_0}
 \leq |Q_0|^{1/p}|f_{Q_0}|
      +|Q_0|^{1/p}\avg_{Q_0}|f-f_{Q_0}|
 \leq |Q_0|^{1/p}|f_{Q_0}|+[f]_{\JN}.
\]
Combining the two bounds proves the claim.
\end{proof}

\begin{lemma}[Full-norm maximal estimate]\label{lem:maximal}
The local dyadic maximal operator $\M$ is bounded on $(\JN,\|\cdot\|_X)$.
\end{lemma}

\begin{proof}
Kinnunen--Myyryl\"ainen, Theorem 4.2, gives
\begin{equation}\label{eq:sourcebound}
 [\M f]_{\JN}\leq C_{n,p}[f]_{\JN}.
\end{equation}
The dyadic John--Nirenberg inequality gives the standard embedding
\begin{equation}\label{eq:weakembed}
 \|f-f_{Q_0}\|_{L^{p,\infty}(Q_0)}\leq C_{n,p}[f]_{\JN}.
\end{equation}
Fix $1<q<p$.  The finite-measure Lorentz embedding and the strong $(q,q)$
maximal inequality imply
\begin{align*}
 \avg_{Q_0}\M(f-f_{Q_0})
 &\leq |Q_0|^{-1/q}\|\M(f-f_{Q_0})\|_{L^q(Q_0)}\\
 &\leq C_{n,q}|Q_0|^{-1/q}\|f-f_{Q_0}\|_{L^q(Q_0)}
 \leq C_{n,p}|Q_0|^{-1/p}[f]_{\JN}.
\end{align*}
Since $\M f\leq\M(f-f_{Q_0})+|f_{Q_0}|$, this controls the average term in
\eqref{eq:Xnorm}.  Together with \eqref{eq:sourcebound}, it proves the
lemma.  Completeness follows from completeness modulo constants proved in the
source paper plus the scalar mean coordinate.
\end{proof}

\section{Rubio de Francia iteration and the decomposition}

\begin{proof}[Proof of Theorem \ref{thm:main}]
Let $A\geq1$ be the operator norm from Lemma \ref{lem:maximal}.  For
$g\geq0$ in $X$, define
\begin{equation}\label{eq:rubio}
 \Rop g:=\sum_{j=0}^{\infty}\frac{(\M)^j g}{(2A)^j}.
\end{equation}
The series converges in $X$, because its $j$th term has norm at most
$2^{-j}\|g\|_X$.  The embedding $X\hookrightarrow L^1(Q_0)$ and monotone
convergence identify this norm limit with the pointwise nonnegative series.
Consequently
\begin{equation}\label{eq:rubioproperties}
 g\leq\Rop g,
 \qquad \|\Rop g\|_X\leq2\|g\|_X,
 \qquad \M(\Rop g)\leq2A\Rop g.
\end{equation}
The last inequality follows from sublinearity, first for partial sums and then
by monotone convergence.  Thus every nonzero $g$ has a positive dyadic $A_1$
majorant with uniformly controlled constant.

For $f\neq0$, apply this to $g=|f|$ and put $w=\Rop|f|$.  Set
\[
 u=2w+f,
 \qquad v=2w.
\]
Then $f=u-v$ and, because $w\geq|f|$,
\[
 w\leq u\leq3w,
 \qquad v=2w.
\]
Hence $u,v>0$ a.e. and
\[
 \M u\leq3\M w\leq6Aw\leq6Au,
 \qquad \M v\leq2A v.
\]
Thus $u,v\in\Aone$.  They lie in $\JN$ by linearity, and Lemma
\ref{lem:absolute} together with \eqref{eq:rubioproperties} gives the claimed
norm bound.  If $f=0$, take $u=v=1$.  The reverse inclusion in
\eqref{eq:characterization} and its norm inequality follow immediately from
linearity and the triangle inequality.

It remains to prove the maximal-power refinement.  The $A_1$ constants of
$u,v$ are bounded solely in terms of $n,p$.  The dyadic reverse H\"older
theorem therefore supplies one uniform $r>1$ and $B<\infty$ such that, for
$z=u,v$ and every $Q\in\D(Q_0)$,
\begin{equation}\label{eq:rh}
 \left(\avg_Qz^r\right)^{1/r}\leq B\avg_Qz.
\end{equation}
In particular $h=z^r\in L^1(Q_0)$.  Dyadic differentiation and
\eqref{eq:rh} yield, almost everywhere,
\[
 z\leq(\M h)^{1/r}
 \leq B\M z
 \leq B[z]_{\Aone}z.
\]
Therefore
\[
 z=k(\M(z^r))^{1/r},
 \qquad k:=\frac{z}{(\M(z^r))^{1/r}},
 \qquad (B[z]_{\Aone})^{-1}\leq k\leq1.
\]
Apply this to $u$ and $v$ to obtain \eqref{eq:maximalrepresentation}.
\end{proof}

\section{Scope and novelty check}

The theorem is an exact difference-cone characterization and the refinement
is an exact maximal-power formula with bounded positive multipliers.  It does
not assert the stronger identity obtained by requiring $k_+=k_-=1$ and the
fixed exponent $1/p$.  The source question does not prescribe that stronger
form.

Focused searches through 13 August 2026 found later work on dyadic $JN_p$
preduals, vanishing subspaces, and fractional dyadic maximal operators, but no
statement of \eqref{eq:characterization} or \eqref{eq:maximalrepresentation}.
The proof is also consistent with the classical Coifman--Rochberg philosophy:
maximal boundedness produces an $A_1$ cone, and the ambient oscillation space
is its difference hull.

\begin{thebibliography}{9}

\bibitem{KM}
J.~Kinnunen and K.~Myyryl\"ainen,
\emph{Dyadic John--Nirenberg space},
Proc. Roy. Soc. Edinburgh Sect. A \textbf{153} (2023), 1--18;
arXiv:2107.00492.

\bibitem{CR}
R.~R. Coifman and R.~Rochberg,
\emph{Another characterization of BMO},
Proc. Amer. Math. Soc. \textbf{79} (1980), 249--254.

\bibitem{DHKoY}
G.~Dafni, T.~Hyt\"onen, R.~Korte, and H.~Yue,
\emph{The space $JN_p$: nontriviality and duality},
J. Funct. Anal. \textbf{275} (2018), 577--603.

\bibitem{RdF}
J.~L. Rubio de Francia,
\emph{Factorization theory and $A_p$ weights},
Amer. J. Math. \textbf{106} (1984), 533--547.

\end{thebibliography}

\end{document}
